\section{Lying over and going up}
We now interpret integrality in terms of the geometry of $\spec$.
In general, for $R \to S$ a ring-homomorphism, the induced map $\spec S \to
\spec R$ need not be topologically nice; for instance, even if $S$ is a
finitely generated $R$-algebra, the image of $\spec S$ in $\spec R$ need not
be either open or closed.\footnote{It is, however, true that if $R$ is
\emph{noetherian} (see \rref{noetherian}) and $S$ finitely generated over
$R$, then the image of $\spec S$ is \emph{constructible,} that is, a finite
union of locally closed subsets. this is Chevalley's constructibility theorem; its proof lies outside this selected supplement.}

We shall see that under conditions of integrality, more can be said.



\subsection{Lying over}

In general, given a morphism of algebraic varieties $f: X \to Y$, the image of
a closed subset $Z \subset X$ is far from closed.  For instance, a regular function $f: X \to \mathbb{C}$ 
that is a closed map would have to be either surjective or constant (if $X$ is
connected, say).
Nonetheless, under integrality hypotheses, we can say more.

\begin{proposition}[Lying over] \label{lyingover}
If $\phi: R \to R'$ is an integral morphism, then the induced map
\[  \spec R' \to \spec R  \]
is a closed map; it is surjective if $\phi$ is injective.
 \end{proposition}

Another way to state the last claim, without mentioning $\spec R'$, is the
following. Assume $\phi$ is injective and integral. Then if
$\mathfrak{p} \subset R$ is prime, then there exists $\mathfrak{q} \subset R'$
such that $\mathfrak{p}$ is the inverse image $\phi^{-1}(\mathfrak{q})$.

\begin{proof} First suppose $\phi$ injective, in which case we must prove the
map $\spec R' \to \spec R$ surjective.
Let us reduce to the case of a local ring.
For a prime $\mathfrak{p} \in \spec R$, we must show that $\mathfrak{p}$
arises as the inverse image of an element of $\spec R'$.
So we replace $R$ with $R_{\mathfrak{p}}$.  We get a map
\[ \phi_{\mathfrak{p}}: R_{\mathfrak{p}} \to (R- \mathfrak{p})^{-1} R'  \]
which is injective if $\phi$ is, since localization is an exact functor. Here
we
have localized both $R, R'$ at the multiplicative subset $R - \mathfrak{p}$.

Note that $\phi_{\mathfrak{p}}$ is an integral extension too. This follows
because integrality is preserved by base-change.
We will now prove the result for $\phi_{\mathfrak{p}}$; in particular, we
will show
that there is a prime ideal of $(R- \mathfrak{p})^{-1} R'$ that pulls back to
$\mathfrak{p}R_{\mathfrak{p}}$. These will imply that if we pull this prime
ideal back to $R'$, it will pull back to $\mathfrak{p}$ in $R$. In detail, we
can consider the diagram
\[ \xymatrix{
\spec (R-\mathfrak{p})^{-1} R'\ar[d]  \ar[r] & \spec R_{\mathfrak{p}} \ar[d] \\
\spec R' \ar[r] &  \spec R
}\]
which shows that if $\mathfrak{p} R_{\mathfrak{p}}$ appears in the image of the top map, then
$\mathfrak{p}$
arises as the image of something in $\spec R'$.
So it is sufficient for the proposition (that is, the case of $\phi$
injective) to handle the case of $R$ local, and
$\mathfrak{p}$ the maximal ideal.


In other words, we need to show that:
\begin{quote}
If $R$ is a \emph{local} ring, $\phi: R \hookrightarrow R'$ an injective
integral morphism, then the maximal ideal of $R$ is the inverse image of
something in $\spec R'$.
\end{quote}

Assume $R$ is local with maximal ideal $\mathfrak{p}$. We want to find a prime ideal $\mathfrak{q} \subset R'$ such that
$\mathfrak{p}  = \phi^{-1}(\mathfrak{q})$. Since $\mathfrak{p}$ is already
maximal, it will suffice to show that $\mathfrak{p} \subset
\phi^{-1}(\mathfrak{q})$. In particular, we need to show that there is a prime
ideal $\mathfrak{q}$ such that
\( \mathfrak{p} R' \subset \mathfrak{q}.  \)
The pull-back of this will be $\mathfrak{p}$.

If $\mathfrak{p}R' \neq R'$, then
$\mathfrak{q}$ exists, since every proper ideal of a ring is contained in a
maximal ideal.  We will in fact show
\begin{equation} \label{thingdoesn'tgenerate} \mathfrak{p} R' \neq R',
\end{equation}
or that $\mathfrak{p}$ does not generate the unit ideal in $R'$.
If we prove \eqref{thingdoesn'tgenerate}, we will thus be able to find our
$\mathfrak{q}$, and we will be done.

Suppose the
contrary, i.e. $\mathfrak{p}R' = R'$. We will derive a contradiction using
Nakayama's lemma (\cref{nakayama}). 
Right now, we cannot apply Nakayama's lemma directly because $R'$ is not a
finite $R$-module. The idea is that we will ``descend'' the ``evidence'' that 
\eqref{thingdoesn'tgenerate} fails to a small subalgebra of $R'$, and then
obtain a contradiction.
To do this, note that $1 \in \mathfrak{p}R'$, and we can write
\[ 1 = \sum x_i \phi(y_i)  \]
where $x_i \in R', y_i \in \mathfrak{p}$.
This is the ``evidence'' that \eqref{thingdoesn'tgenerate} fails, and it
involves only a finite amount of data. 

Let $R''$ be the subalgebra of $R'$ generated by $\phi(R)$ and the $x_i$. Then
$R'' \subset R'$ and is finitely generated as an $R$-algebra, because it is generated by
the $x_i$. However, $R''$ is integral over $R$ and thus finitely generated as an
$R$-module, by \cref{finitelygeneratedintegral}. This
is where integrality comes in.

So $R''$ is a finitely generated $R$-module. Also, the expression
$1 = \sum x_i \phi(y_i)$ shows that $\mathfrak{p}R'' = R''$. However, this
contradicts Nakayama's lemma. That brings the contradiction, showing that
$\mathfrak{p}$ cannot generate $(1)$ in $R'$, and proving the surjectivity
part of  lying over theorem.

Finally, we need to show that if $\phi: R \to R'$ is \emph{any} integral
morphism, then $\spec R' \to \spec R$ is a closed map. Let $X = V(I) $ be a
closed subset of $\spec R'$. Then the image of $X$ in $\spec R$ is the image
of the map
\[ \spec R'/I \to \spec R   \]
obtained from the morphism $R \to R' \to  R'/I$, which is integral; thus we are
reduced to showing that any integral morphism $\phi$ has closed image on the
$\spec$. 
Thus we are reduced to $X = \spec R'$, if we throw out $R'$ and replace it by
$R'/I$.

In other words, we must prove the following statement. Let $\phi: R \to R'$ be
an integral morphism; then the image of $\spec R'$ in $\spec R$ is closed.
But, quotienting by $\ker \phi$ and taking the map $R/\ker \phi \to R'$, we
may reduce to the case of $\phi$ injective; however, then this follows from
the surjectivity result already proved.
\end{proof}

In general, there will be \emph{many} lifts of a given prime ideal.
Consider for instance the inclusion $\mathbb{Z} \subset \mathbb{Z}[i]$.
Then the prime ideal $(5) \in \spec \mathbb{Z}$ can be lifted either to $(2+i)
\in \spec \mathbb{Z}[i]$
or $(2-i) \in \spec \mathbb{Z}[i]$.
These are distinct prime ideals: $\frac{2+i}{2-i} \notin \mathbb{Z}[i]$. But
note that any element of $\mathbb{Z}$ divisible by $2+i$ is automatically
divisible by its conjugate $2-i$, and consequently by their product $5$
(because $\mathbb{Z}[i]$ is  a UFD, being a euclidean domain).

Nonetheless, the different lifts are incomparable.

\begin{proposition} \label{incomparableintegral}
Let $\phi: R \to R'$ be an integral morphism. Then given $\mathfrak{p} \in
\spec R$, there are no inclusions among the elements $\mathfrak{q} \in \spec
R'$ lifting $\mathfrak{p}$.
\end{proposition} 
In other words, if $\mathfrak{q}, \mathfrak{q}' \in \spec R'$ lift
$\mathfrak{p}$, then $\mathfrak{q} \not\subset \mathfrak{q}'$.
\begin{proof} 
We will give a ``slick'' proof by various reductions.
Note that the operations of localization and quotienting only shrink the
$\spec$'s: they do not ``merge'' heretofore distinct prime ideals into one.
Thus, by quotienting $R$ by $\mathfrak{p}$, we may assume $R$ is a
\emph{domain} and that $\mathfrak{p} = 0$.
Suppose we had two primes $\mathfrak{q} \subsetneq \mathfrak{q}'$ of $R'$
lifting $(0) \in \spec R$.
Quotienting $R'$ by $\mathfrak{q}$, we may assume that $\mathfrak{q}  = 0$.
We could even assume $R \subset R'$, by quotienting by the kernel of $\phi$.
The next lemma thus completes the proof, because it shows that $\mathfrak{q}'
= 0$, contradiction.
\end{proof} 

\begin{lemma} 
Let $R \subset R'$ be an inclusion of integral domains, which is an integral
morphism. If $\mathfrak{q} \in \spec R'$ is a nonzero prime ideal, then
$\mathfrak{q} \cap R$ is nonzero.
\end{lemma} 
\begin{proof} 
Let $x \in \mathfrak{q}'$ be nonzero. There is an equation
\[ x^n + r_1 x^{n-1} + \dots + r_n = 0,  \quad r_i \in R, \]
that $x$ satisfies, by assumption. 
Here we can assume $r_n \neq 0$; then $r_n \in \mathfrak{q}' \cap R$ by
inspection, though. So this intersection is nonzero.

\end{proof} 

\begin{corollary} 
Let $R \subset R'$ be an inclusion of integral domains, such that $R'$ is
integral over $R$. Then if one of $R, R'$ is a field, so is the other.
\end{corollary} 
\begin{proof} 
Indeed, $\spec R' \to \spec R$ is surjective by \cref{lyingover}: so if $\spec
R'$ has one element (i.e., $R'$ is a field), the same holds for $\spec R$
(i.e., $R$ is a field). Conversely, suppose $R$ a field. 
Then any two prime ideals in $\spec R'$ pull back to the same element of $\spec
R$. So, by \cref{incomparableintegral}, there can be no inclusions among the prime ideals of $\spec
R'$. But $R'$ is a domain, so it must then be a field.
\end{proof} 


\begin{exercise} Let $k$ be a field.
Show that $k[\mathbb{Q}_{\geq 0}]$ is integral over the polynomial ring $k[T]$.
Although this is a \emph{huge} extension,  the prime ideal $(T)$ lifts in only
one way to $\spec k[\mathbb{Q}_{\geq 0}]$.
\end{exercise} 

\begin{exercise} 
Suppose $A \subset B$ is an inclusion of rings over a field of characteristic
$p$. Suppose $B^p \subset A$, so that $B/A$ is integral in a very strong sense. 
Show that the map $\spec B \to \spec A$ is a \emph{homeomorphism.} 
\end{exercise} 



\subsection{Going up}
Let $R \subset R'$ be an inclusion of rings with $R'$ integral over $R$. We saw in the
lying over theorem (\cref{lyingover}) that any prime $\mathfrak{p} \in \spec R$ 
has a prime $\mathfrak{q} \in \spec R'$ ``lying over'' $\mathfrak{p}$, i.e.
such that $R \cap \mathfrak{q} = \mathfrak{p}$.
We now want to show that we can lift finite \emph{inclusions} of primes to $R'$.

\begin{proposition}[Going up]
Let $R \subset R'$ be an integral inclusion of rings.
Suppose $\mathfrak{p}_1 \subset \mathfrak{p}_2 \subset \dots \subset
\mathfrak{p}_n \subset R$ is a 
finite ascending chain of prime ideals in $R$.
Then there is an ascending chain $\mathfrak{q}_1 \subset \mathfrak{q}_2 \subset
\dots \subset \mathfrak{q}_n$ in $\spec R'$ lifting this chain. 

Moreover, $\mathfrak{q}_1$ can be chosen arbitrarily so as to lift
$\mathfrak{p}_1$.
\end{proposition} 
\begin{proof} 
By induction and lying over (\cref{lyingover}), it suffices to show:
\begin{quote}
Let $\mathfrak{p}_1 \subset \mathfrak{p}_2$ be an inclusion of primes in $\spec
R$. Let $\mathfrak{q}_1 \in \spec R'$ lift $\mathfrak{p}_1$. Then there is
$\mathfrak{q}_2 \in \spec R'$, which satisfies the dual conditions of lifting
$\mathfrak{p}_2$ and containing $\mathfrak{q}_1$.
\end{quote}
To show that this is true, we apply \cref{lyingover} to the inclusion
$R/\mathfrak{p}_1 \hookrightarrow R'/\mathfrak{q}_1$. There is an element of
$\spec R'/\mathfrak{q}_1$ lifting $\mathfrak{p}_2/\mathfrak{p}_1$; the
corresponding element of $\spec R'$ will do for $\mathfrak{q}_2$.
\end{proof} 

